\documentclass{beamer}
\usetheme{Madrid}
\usepackage{amsmath, amssymb, amsthm}
\usepackage{graphicx}
\usepackage{listings}
\usepackage{gensymb}
\usepackage[utf8]{inputenc}
\usepackage{hyperref}
\usepackage{gvv}
\usepackage{tikz}
\lstset{
  language=Python,
  basicstyle=\ttfamily\small,
  keywordstyle=\color{blue},
  stringstyle=\color{orange},
  numbers=left,
  numberstyle=\tiny\color{gray},
  breaklines=true,
  showstringspaces=false
}
\usetikzlibrary{decorations.pathmorphing}

\title{Question-4.13.27}
\author{EE25BTECH11022 - sankeerthan}
\date{}
\begin{document}

\frame{\titlepage}

\begin{frame}
\frametitle{Question}
Consider the set of all lines $px + qy + r = 0$ such that $3p + 2q + 4r = 0$.
Which one of the following statements is true?
\begin{enumerate}
  \item The lines are concurrent at the point $\brak{\frac{3}{4},\frac{1}{2}}$.
  \item Each line passes through the origin.
  \item The lines are all parallel.
  \item The lines are not concurrent.
\end{enumerate}
\end{frame}

\begin{frame}
\frametitle{Solution}
Given the set of lines $px + qy + r = 0$ such that $3p + 2q + 4r = 0$.
\begin{align}
\myvec{x & y & 1 \\3 & 2 & 4}\myvec{p \\ q \\ r}=\myvec{0 \\ 0}
\end{align}
For a non-trivial solution,the vector $\brak{p, q, r}$ must be in the nullspace. The system has a non-trivial solution if and only if the row vectors are linearly dependent:
\begin{align} 
\brak{x_0, y_0, 1} = \lambda\brak{3, 2, 4} \\
x_0 = 3\lambda,y_0 = 2\lambda,1 = 4\lambda  \\
1 = 4\lambda \implies \lambda = \frac{1}{4} \\
x_0 = \frac{3}{4},\qquad y_0 = \frac{1}{2} 
\end{align}
Therefore, all lines in the given family are concurrent at the point
$\brak{\frac{3}{4},\frac{1}{2}}$
\end{frame}

\begin{frame}
\frametitle{graph}
\begin{figure}
    \centering
    \includegraphics[width=0.7\columnwidth]{figs/points.png}
    \caption{}
\end{figure}
\end{frame}
\begin{frame}[fragile]
\frametitle{C-Code}
\begin{lstlisting}[language=C]
// code.c
#include <stdio.h>

// Family of lines: parametrize p=t, q=1, r = -(3t + 2)/4.
// Generates n points on each such line passing through concurrency point (3/4, 1/2).
void generate_points(double *x, double *y, int n, double t) {
    double q = 1.0;
    double r = -(3.0*t + 2.0)/4.0;
    double x_start = -1.0;
    double x_end = 2.0;
    for (int i = 0; i < n; i++) {
        x[i] = x_start + (x_end - x_start)*i/(n-1);
        y[i] = (-t*x[i]-r)/q;
    }
}
\end{lstlisting}
\end{frame}
\begin{frame}[fragile]
\frametitle{Python-Code}
\begin{lstlisting}
# Code by GVV Sharma
# Modified for Problem Solution
# Released under GNU GPL
# Calculating area enclosed between curves
# call.py
import ctypes
import numpy as np

lib = ctypes.CDLL('./code.so')
\end{lstlisting}
\end{frame}
\begin{frame}[fragile]
\frametitle{Python-Code}
\begin{lstlisting}
lib.generate_points.argtypes = [
    np.ctypeslib.ndpointer(dtype=np.float64, flags="C_CONTIGUOUS"),
    np.ctypeslib.ndpointer(dtype=np.float64, flags="C_CONTIGUOUS"),
    ctypes.c_int,
    ctypes.c_double
]
def extract_points(n=10, t=2.0):
    x = np.zeros(n, dtype=np.float64)
    y = np.zeros(n, dtype=np.float64)
    lib.generate_points(x, y, n, t)
    return x, y

if __name__ == "__main__":
    x, y = extract_points(10, 2.0)
    print("x:", x)
    print("y:", y)
\end{lstlisting}
\end{frame}
\begin{frame}[fragile]
\frametitle{Python-Code}
\begin{lstlisting}
# plot.py
import matplotlib.pyplot as plt
from call import extract_points

n_points = 20
concurrency_point = (0.75, 0.5)
t_values = [1.0, 1.5, 2.0, 2.5]  # Different lines in the family

plt.figure(figsize=(8,6))

for t in t_values:
    x, y = extract_points(n_points, t)
    plt.plot(x, y, label=f'p={t}')
    plt.scatter(x, y, s=10)
\end{lstlisting}
\end{frame}
\begin{frame}[fragile]
\frametitle{Python-Code}
\begin{lstlisting}

plt.scatter([concurrency_point[0]], [concurrency_point[1]], color='red', marker='X', s=100, label='Concurrency point')
plt.title('Family of lines passing through (3/4, 1/2)')
plt.xlabel('x')
plt.ylabel('y')
plt.grid(True)
plt.legend()
plt.show()
\end{lstlisting}
\end{frame}

\end{document}